%\section{Protocols and Institutional Designs Compatible with FRFP}
This appendix describes protocols and institutional patterns that are compatible with the
formal architecture developed in the main text. All constructions here are grounded in
the existing FRFP mathematics: explicit--tacit separation, pipeline semantics in
\(N \ltimes E\), correctness invariance under navigation, the dynamic analysis of
hallucination, and the impossibility of fully automating correctness from explicit
artifacts alone. We do not introduce new theorems; instead, we translate the formal
constraints into practical rules for long-horizon human--AI collaboration.

\subsection{Foundational Rules for FRFP-Compatible Practice}

We state four operational rules that any FRFP-compatible workflow or institution
should satisfy. Each rule is a protocol-level reflection of a mathematical property
of the architecture.

\paragraph{Rule 1 (Explicit confinement).}
All machine-mediated operations occur in the explicit category \(E\): symbolic
transformations, derivations, simulations, code execution, and navigation in
\(N \ltimes E\). The system never claims tacit authority or correctness. The only
boundary from explicit to tacit space is the backflow map \(RB : E \to T\), which
hands explicit artifacts to human agents for tacit evaluation.

\paragraph{Rule 2 (Human-exclusive grounding).}
Correctness, acceptance, and responsibility live exclusively in tacit space \(T\).
Observable judgments are given by \(\gamma \circ \llbracket - \rrbracket\), where
\(\llbracket - \rrbracket : N \ltimes E \to T\) is the semantic functor and
\(\gamma : \mathrm{Ob}(T) \to J\) is the correctness quotient. In protocol terms:
only humans may accept, reject, or suspend judgment on claims; AI systems present
explicit proposals and diagnostics but never act as final arbiters.

\paragraph{Rule 3 (Source–reason traceability).}
Every published claim must be traceable to:
(i) a normal-form explicit pipeline \(P \in N \ltimes E\), and
(ii) a sequence of tacit evaluations triggered by its boundary step \(RB\).
Formally, reproducible provenance means that independent investigators can reconstruct
a pipeline whose normal form has the same explicit semantics and yields the same
observable tacit judgment \(\gamma(\llbracket P \rrbracket)\). In practice this
requires versioned logs of prompts, model configurations, code, and intermediate
artifacts, together with human evaluation checkpoints.

\paragraph{Rule 4 (Integrity under navigation).}
Because FRFP guarantees that \(\gamma \circ \llbracket - \rrbracket\) is invariant
under explicit reduction and navigation in \(N \ltimes E\), protocols must preserve
this invariance at the interface level: branching, backtracking, and exploratory
reformulations are allowed and encouraged, but they are always treated as alternative
paths toward a single tacit evaluation, not as independent correctness claims. User
interfaces and institutional processes should therefore separate:
(i) exploratory branches in explicit space, and
(ii) the consolidated tacit judgments to which they are ultimately submitted.

\subsection{Dialogic Reproducibility Protocol}

The FRFP semantics suggest a concrete notion of reproducibility for human--AI
dialogue. A dialogic workflow is reproducible when an independent team can reconstruct
an explicit pipeline \(P' \in N \ltimes E\) whose normal form is equivalent to the
original pipeline \(P\) up to navigation, and for which the resulting tacit
judgments coincide:
\[
  \gamma(\llbracket P' \rrbracket) = \gamma(\llbracket P \rrbracket).
\]

A minimal protocol compatible with this notion has three components.

\paragraph{(1) Structured interaction traces.}
Each interaction with the AI system is recorded as a typed step in a TDG-generated
pipeline: representation initiation (RI), explicit computation (EC), explicit
diagnostics (ED), navigation steps in \(N\), and boundary steps \(RB\). Logs include
model version, code or tool calls, and the explicit artifacts produced at each stage.

\paragraph{(2) Context archives.}
At designated checkpoints, the current explicit state (definitions, assumptions,
intermediate results) is serialized as a context artifact in \(E\). These context
snapshots serve as canonical entry points for future pipelines and allow independent
teams to replay or extend the reasoning using the same initial explicit data.

\paragraph{(3) Replay procedures.}
A reproducibility audit consists of reconstructing a pipeline from archived context
and logs, reducing it to normal form in \(N \ltimes E\), and checking whether the
resulting tacit judgment matches the recorded outcome. Discrepancies expose either
implementation errors, insufficient logging, or tacit disagreements, each of which
can then be investigated explicitly.

\subsection{Closure Criteria for Long-Horizon Collaboration}

Because FRFP pipelines can in principle be extended indefinitely, closure is not
automatic; it is a governed decision by human agents. We define closure criteria
that align with the FRFP semantics and the correctness invariance theorem.

\paragraph{Conceptual stability.}
The explicit normal forms associated with the main claims remain stable across
multiple iterations of exploration and navigation. Operationally, further TDG
expansions and navigation steps in \(N \ltimes E\) no longer change the canonical
explicit structures relevant to the claims. This reflects the fact that
\(\gamma \circ \llbracket - \rrbracket\) depends only on the semantic content of the
pipeline, not on its construction history.

\paragraph{Evidential sufficiency.}
Every major claim is linked to a well-typed explicit pipeline whose semantic image
\(\llbracket P \rrbracket \in T\) justifies the corresponding tacit judgment under
the agreed standards of the field (proof obligations, error bounds, empirical
thresholds). In practice, this means that each claim is backed by explicit
derivations, simulations, or data analyses that can be inspected and replayed.

\paragraph{Provenance completeness.}
For each published claim, the associated pipeline and its key checkpoints are
recorded in sufficient detail to satisfy the dialogic reproducibility protocol
above. This ensures that future investigators can reconstruct \(P\), check its
normal form, and compare their tacit judgments with those originally recorded.

Closure is declared when these three conditions hold for the main claims and no
new FRFP-compatible pipeline, starting from the archived context, is expected to
change the relevant tacit judgments. At that point, the interaction transitions
from open-ended exploration to a stable record suitable for publication and
institutional review.

\subsection{Illustrative Case Study: Stress-Testing a Financial Strategy}

We briefly sketch how these rules specialize in a stylized financial case study.
The goal is not to introduce new mathematics but to show how FRFP-compatible
protocols can be instantiated in a concrete domain.

\paragraph{Scenario.}
A team is developing a portfolio rebalancing strategy under hierarchical limits
and wishes to stress-test it using both human expertise and AI tools. The explicit
artifacts \(E\) include model code, parameter sets, shock scenarios, and summary
statistics; tacit space \(T\) represents human judgments about adequacy, risk, and
regulatory acceptability.

\paragraph{Pipelines and navigation.}
The team uses the AI system to generate candidate stress-test pipelines: baseline
scenarios, moderate shocks, and tail events (e.g., adding a tail scenario with a
spread shock at or above the \(99\)th percentile in a sector). Each pipeline is a
TDG term in \(N \ltimes E\), built from explicit primitives (scenario generation,
simulation runs, aggregation) and navigation steps (branching on alternative models,
rolling back to prior assumptions, refining constraints).

\paragraph{Dialogic reproducibility.}
At each major checkpoint, the team archives:
(i) the explicit configuration (portfolio state, limits, scenarios);
(ii) the simulation outputs; and
(iii) the corresponding AI-generated analyses.
Human evaluators record tacit judgments on whether the current set of scenarios,
including tail shocks, sufficiently probes the relevant risk surfaces. A reproducible
record thus consists of:
\begin{itemize}
  \item the explicit pipeline \(P\) with its normal form (the set of scenarios and
        computations actually used), and
  \item the sequence of tacit judgments \(\gamma(\llbracket P \rrbracket)\) at closure.
\end{itemize}

\paragraph{Closure and institutional use.}
Closure is reached when:
(i) additional navigation moves (e.g., adding more variations of the \(99\)th percentile
shocks) no longer change the team’s tacit assessment of the strategy’s robustness;
(ii) each risk statement (e.g., “the strategy remains feasible under shocks up to
X”) is backed by explicit simulations; and
(iii) all relevant pipelines and checkpoints are logged for replay.
An institution (such as a risk committee) can then review the explicit pipelines
and associated tacit judgments, but by the non-automation results in the FRFP
appendices it cannot replace human evaluation with a purely explicit surrogate
without losing the correctness guarantees.

\medskip

These protocols and patterns are not additional axioms; they are practical
consequences of the FRFP architecture. They show how explicit confinement, human-
exclusive grounding, correctness invariance, and the structure of hallucination
and institutional limits can be used to design concrete workflows and institutions
for long-horizon human--AI collaboration.
